\section*{Main guarantee and conventions (original Sections 1--2)}
A code $\mathcal C\subseteq\{0,1\}^{\bar n}$ is $(s,L)$ erasure list-decodable if for every word $w$, after erasing any $s$ symbols of $w$, the remaining $\bar n-s$ symbols have at most $L$ possible completions into a codeword of $\mathcal C$.
\setcounter{theorem}{1}
\begin{theorem}[see also Theorem 46]
For every constant $0<\gamma<1$ and $\varepsilon=n^{-\Omega_\gamma(1)}$ there exists an explicit $((1-\varepsilon)\bar n,L=\log^{O_\gamma(1)}\frac1\varepsilon)$ erasure list-decodable code $\mathcal C:\{0,1\}^n\to\{0,1\}^{\bar n}$ of rate $\varepsilon^{1+\gamma}$.
\end{theorem}
We stress that the codes we present are explicit in the sense that they have explicit encoding, but we do not know whether the codes we construct admit efficient erasure list-decoding algorithms. We also mention that the list-size $\operatorname{poly}(\log\frac1\varepsilon)$ achieved by our code is exponentially smaller than the best possible list-size by any linear code.
Throughout the paper we use the convention that lowercase variables are the logarithm (in base 2) of their corresponding uppercase variables, e.g., $n=\log N$, $d=\log D$. We denote by $[A]$ the set $\{1,\ldots,A\}$. The density of a set $B\subseteq A$ is $\rho(B)=|B|/|A|$. We say a function $f:A\to B$ is explicit if there exists a deterministic polynomial algorithm that runs in time $\operatorname{poly}(\log|A|)$ and computes $f$.
\section*{Preliminaries (original Section 2)}
\subsection*{Random variables and min-entropy}
The statistical distance between two distributions $X$ and $Y$ on the same domain $\Omega$ is defined as $|X-Y|=\max_{A\subseteq\Omega}(\Pr[X\in A]-\Pr[Y\in A])$. If $|X-Y|\le\varepsilon$ we say $X$ is $\varepsilon$-close to $Y$ and denote it by $X\approx_\varepsilon Y$. We denote by $U_n$ the random variable distributed uniformly over $\{0,1\}^n$. We say a random variable is flat if it is uniform over its support.
For a function $f:\Omega_1\to\Omega_2$ and a random variable $X$ distributed over $\Omega_1$, $f(X)$ is the random variable distributed over $\Omega_2$ obtained by choosing $x$ according to $X$ and computing $f(x)$. For a set $A\subseteq\Omega_1$, $f(A)=\{f(x)\mid x\in A\}$. For every $f:\Omega_1\to\Omega_2$ and two random variables $X$ and $Y$ distributed over $\Omega_1$, it holds that $|f(X)-f(Y)|\le|X-Y|$.
The min-entropy of a random variable $X$ is defined by
\[H_\infty(X)=\min_{x\in\Supp(X)}\log\frac1{\Pr[X=x]}.\]
A random variable $X$ is an $(n,k)$ source if $X$ is distributed over $\{0,1\}^n$ and has min-entropy at least $k$. When $n$ is clear from the context we sometimes omit it and simply say that $X$ is a $k$-source. Every $k$-source $X$ can be expressed as a convex combination of flat distributions each with min-entropy at least $k$.
\setcounter{theorem}{5}
\begin{definition}[average conditional min-entropy]
Let $X,Y$ be two random variables. The average conditional min-entropy of $X$ given $Y$ is
\[\HH(X\mid Y)=-\log\left(\E_{y\sim Y}\left[2^{-H_\infty(X\mid Y=y)}\right]\right).\]
\end{definition}
We will use the following simple claim about average conditional min-entropy:
\begin{claim}For any random variables $X,Y$, $\HH(X\mid Y)\ge H_\infty(X)-\log|\Supp(Y)|$.\end{claim}
\subsection*{Condensers}
\begin{definition}[condenser]
A function $C:\{0,1\}^n\times\{0,1\}^d\to\{0,1\}^m$ is an $(n,k)\to_{\varepsilon_{\rm cond}}(m,k')$ condenser, if for every $(n,k)$ source $X$, $C(X,U_d)$ is $\varepsilon_{\rm cond}$-close to an $(m,k')$ source. If $k=\delta n$ and $k'=\delta'm$ we say $C$ is a $\delta\to_{\varepsilon_{\rm cond}}\delta'$ condenser.
\end{definition}
\begin{lemma}
Suppose $C:\{0,1\}^n\times\{0,1\}^d\to\{0,1\}^m$ is an $(n,k)\to_{\varepsilon_{\rm cond}}(m,k'+d)$ condenser. Let $X$ be an $(n,k)$ source. Let $\varepsilon_i$ be the minimal distance of $C(X,i)$ to an $(m,k')$ source. Then, $\E_{i\in\{0,1\}^d}[\varepsilon_i]\le\varepsilon_{\rm cond}$.
\end{lemma}
\begin{proof}
Fix an $(n,k)$ source $X$. For $i\in\{0,1\}^d$, let $H_i\subseteq\{0,1\}^m$ be the set of ``heavy'' elements of $C(X,i)$,
\[H_i=\{w\in\{0,1\}^m:\Pr_{x\in X}[C(x,i)=w]\ge2^{-k'}\}.\]
The distance of $C(X,i)$ from a $k'$-source is
\[\varepsilon_i=\Pr_{x\in X}[C(x,i)\in H_i]-2^{-k'}|H_i|,\]
by redistributing the mass of the heavy elements. Let $H=\bigcup_{i\in\{0,1\}^d}H_i$. Then,
\begin{itemize}
\item For every $w\in H$, $\Pr_{x\in X,i\in\{0,1\}^d}[C(x,i)=w]\ge2^{-d}2^{-k'}=2^{-(k'+d)}$, and,
\item it holds that
\begin{align*}
\varepsilon_{\rm cond}&=\Pr_{x\in X,i\in\{0,1\}^d}[C(x,i)\in H]-|H|2^{-(k'+d)}\\
&=\sum_{i\in\{0,1\}^d}2^{-d}\Pr_x[C(x,i)\in H]-|H|2^{-(k'+d)}\\
&\ge\sum_{i\in\{0,1\}^d}2^{-d}\Pr_x[C(x,i)\in H_i]-2^{-(k'+d)}\sum_{i\in\{0,1\}^d}|H_i|\\
&=\sum_{i\in\{0,1\}^d}2^{-d}(\Pr_x[C(x,i)\in H_i]-2^{-k'}|H_i|)\\
&=\sum_{i\in\{0,1\}^d}2^{-d}\varepsilon_i=\E_{i\in\{0,1\}^d}[\varepsilon_i].
\end{align*}
\end{itemize}
\end{proof}
\subsection*{Two-source extractors}
\begin{definition}[two-source extractor]
A function $\twoExt:\{0,1\}^{n_1}\times\{0,1\}^{n_2}\to\{0,1\}^m$ is an $((n_1,k_1),(n_2,k_2),\varepsilon)$ two-source extractor if for every two independent sources $X_1$ and $X_2$ where $X_1$ is an $(n_1,k_1)$ source and $X_2$ is an $(n_2,k_2)$ source, it holds that $\twoExt(X_1,X_2)\approx_\varepsilon U_m$. We say that $\twoExt$ is strong if
\[(\twoExt(X_1,X_2),X_1)\approx_\varepsilon(U_m,X_1)\quad\text{and}\quad(\twoExt(X_1,X_2),X_2)\approx_\varepsilon(U_m,X_2).\]
\end{definition}
In our construction, we will use the following two-source extractor:
\begin{theorem}[{[39]}]
For every constant $\delta_{\rm Raz}>\frac12$ there exist constants $c_1=c_1(\delta_{\rm Raz}),c_2=c_2(\delta_{\rm Raz})>1$ such that for every $n_1,k_1,n_2,k_2$ satisfying
\[k_1\ge c_1\log n_2,\qquad k_2\ge c_2\log n_1,\]
there exists an explicit strong $((n_1,k_1),(n_2,k_2=\delta_{\rm Raz}n_2),\varepsilon_{\rm Raz})$ two-source extractor
\[\Raz:\{0,1\}^{n_1}\times\{0,1\}^{n_2}\to\{0,1\}^m\]
with $m=\Omega(\min\{k_1,k_2\})$ and $\varepsilon_{\rm Raz}=2^{-\Omega(m)}$, where the constants hiding in the asymptotic notation may depend on $\delta_{\rm Raz}$.
\end{theorem}
\begin{claim}
Suppose $\twoExt:\{0,1\}^{n_1}\times\{0,1\}^{n_2}\to\{0,1\}^m$ is a strong $((n_1,k_1),(n_2,k_2),\varepsilon)$ two-source extractor. Let $X$ be an $(n,k_1)$ source. Call an element $y\in\{0,1\}^{n_2}$ bad if $|\twoExt(X,y)-U_m|>\varepsilon$, and let $B_Y$ denote the set of all bad elements. Then, $|B_Y|<2^{k_2}$.
\end{claim}
\begin{proof}
Assume towards contradiction that $|B_Y|\ge2^{k_2}$ and let $Y$ be the uniform distribution over the set $B_Y$. Then, $H_\infty(Y)\ge k_2$ and so $(\twoExt(X,Y),Y)\approx_\varepsilon(U_m,Y)$ which implies that
\[\frac1{|B_Y|}\sum_{y\in B_Y}|\twoExt(X,y)-U_m|\le\varepsilon.\]
However, $|\twoExt(X,y)-U_m|>\varepsilon$ for every $y\in B_Y$, in contradiction.
\end{proof}
\subsection*{Mergers}
A merger takes as input a list of possibly correlated random variables along with a short uniform seed and outputs one random variable which is close to having high min-entropy, provided at least one of the input variables has high min-entropy. Formally:
\begin{definition}[somewhere-random source]
A source $X=X_1\circ\ldots\circ X_A$ is an $(n,k,(\alpha,\beta))$ somewhere-random (s.r.) source if there is a random variable $I\in\{0,\ldots,A\}$ such that for every $i\in[A]$, $(X_i\mid I=i)$ is $\alpha$-close to an $(n,k)$ source and $\Pr[I=0]\le\beta$. The variable $I$ is called the indicator of source. If $\alpha=\beta=0$ we say $X$ is a $(n,k)$ s.r. source.
\end{definition}
\begin{definition}[merger]
A function $B:(\{0,1\}^n)^D\times\{0,1\}^t\to\{0,1\}^m$ is a $(k,k',\varepsilon)$ merger, if for every $(n,k)$ s.r. source $X=X_1\circ\ldots\circ X_A$, the output $M(X,U_t)$ is $\varepsilon$-close to a $k'$-source.
\end{definition}
There are explicit constructions of good mergers. Dvir and Wigderson [17] constructed the curve merger and proved that it works with $t=O(\log\frac n\varepsilon)$. This was further improved in [16] who proved that $t=O(\log\frac D\varepsilon)$ suffices. Notice that now $t$ only depends on the number of sources $D$ and the requested error $\varepsilon$, but not on the source length $n$, and this remarkable property will be crucial for us. Formally,
\begin{theorem}[{[17,16]}]
There exists a constant $c_{\rm DKSS}\ge1$ such that the following holds. Fix $\beta,\delta,\varepsilon>0$. There exists an explicit function $B:(\{0,1\}^n)^D\times\{0,1\}^t\to\{0,1\}^n$ that is a $(k=\delta n,k'=(1-\beta)\delta n,\varepsilon)$ merger, with $t=c_{\rm DKSS}\cdot\frac1\beta\log\frac D\varepsilon$.
\end{theorem}
\subsection*{Correlation breakers with advice}
A correlation-breaker with advice is a function $\CBA:\{0,1\}^n\times\{0,1\}^\ell\times[A]\to\{0,1\}^m$ where we think of the first input as a weak source, the second as an independent short seed and the last as an advice string. Roughly speaking, applying CBA on $t$ possibly correlated seeds with $t$ distinct advice strings results in independent random variables. For example, $\CBA(X,Y,\alpha)$ is (nearly) independent of $\CBA(X,Y,\alpha')$ for any $\alpha\ne\alpha'$. Formally,
\begin{definition}
A function $\CBA:\{0,1\}^n\times\{0,1\}^\ell\times[A]\to\{0,1\}^m$ is a $(t,k,\varepsilon_{\rm CBA})$ correlation-breaker with advice if the following holds. If $Y$ is a distribution over $\{0,1\}^n$, $Z=(Z_1,\ldots,Z_t)$ is a distribution on $(\{0,1\}^\ell)^t$, $H$ is a random variable and $\delta>0$ satisfy:
\begin{itemize}
\item $Y$ and $Z$ are independent, conditioned on $H$,
\item $\HH(Y\mid H)\ge k+\log(1/\varepsilon_{\rm CBA})$,
\item $(Z_1,H)\approx_\delta(U_\ell,H)$, and,
\item $\alpha_1,\ldots,\alpha_t\in[A]$ are distinct strings.
\end{itemize}
Then,
\begin{multline*}
(\CBA(Y,Z_1,\alpha_1),(\CBA(Y,Z_i,\alpha_i))_{i=2}^t,H)\\
\approx_{\delta+2\varepsilon_{\rm CBA}}(U_m,(\CBA(Y,Z_i,\alpha_i))_{i=2}^t,H).
\end{multline*}
\end{definition}
We use the following result:
\begin{theorem}[{[13, Theorem 4.12]}]
There exists a constant $c_{\rm CBA}\ge1$ such that the following holds. Let $n,a$ be integers and $\varepsilon_{\rm CBA}>0$. Then, there exists an explicit $(t,k_{\rm CBA},\varepsilon_{\rm CBA})$ correlation-breaker with advice
\[\CBA:\{0,1\}^n\times\{0,1\}^\ell\times[A]\to\{0,1\}\]
with $\ell=c_{\rm CBA}\cdot at\cdot\log\frac n{\varepsilon_{\rm CBA}}$ and $k_{\rm CBA}\ge\ell$.
\end{theorem}
In our setting, the number of rows $A$ is a constant independent of $n$. For this reason we work with a ``basic'' correlation-breaker, where there is no attempt to optimize the dependence of $\ell$ on $a$. This gives a seed-length which is optimal up to constant multiplicative factors.
We also need the following lemma.
\begin{lemma}
Let $X_1,\ldots,X_t$ be random variables over $\{0,1\}^m$. Further suppose that for any $i\in[t]$,
\[(X_i,\{X_j\}_{j\ne i})\approx_\varepsilon(U_m,\{X_j\}_{j\ne i}).\]
Then, $(X_1,\ldots,X_t)\approx_{t\varepsilon}U_{tm}$.
\end{lemma}
\subsection*{Limited independence and non-oblivious bit-fixing sources}
\begin{definition}
A distribution $X$ over $\{0,1\}^A$ is called $(t,\gamma)$-wise independent if the restriction of $X$ to every $t$ coordinates is $\gamma$-close to $U_t$. A source $X$ over $\{0,1\}^A$ is called a $(q,t,\gamma)$ non-oblivious bit-fixing source if there exists a subset $Q\subseteq A$ of size at most $q$ such that the joint distribution of the bits in $A\setminus Q$ is $(t,\gamma)$-wise independent. The bits in $Q$ are allowed to arbitrarily depend on the bits in $A\setminus Q$. If $\gamma=0$ we often say that $X$ is a $(q,t)$ non-oblivious bit-fixing source.
\end{definition}
\begin{lemma}[{[3]}]A $(t,\gamma)$-wise distribution over $A$ bits is $(A^t\gamma)$-close to some $t$-wise independent distribution.\end{lemma}
\begin{definition}
Let $f:\{0,1\}^A\to\{0,1\}$, $\mathcal D$ a distribution over $\{0,1\}^A$ and $Q\subseteq A$. Let $I_{Q,\mathcal D}(f)$ denote the probability that $f$ is undetermined when the variables outside $Q$ are sampled from $\mathcal D$. We define $I_{q,t,\gamma}(f)$ to be the maximum of $I_{Q,\mathcal D}(f)$ over all $Q\subseteq A$ of size $q$ and all $\mathcal D$ that is a $(t,\gamma)$-wise independent distribution. We say that $f$ is $(t,\gamma)$-independent $(q,\varepsilon)$-resilient if $I_{q,t,\gamma}(f)\le\varepsilon$.
\end{definition}
\begin{theorem}[{[9,32]}]
For every $0<\gamma<1$ there exists a constant $c_\gamma\ge1$ such that for all $A>0$ there exists an explicit function $f:\{0,1\}^A\to\{0,1\}$ with the following property: For every $t\ge c_\gamma\log^4 A$,
\begin{itemize}
\item $f$ is almost balanced: For any $t$-wise independent distribution $\mathcal D$ on $\{0,1\}^A$, $\Pr_{x\sim\mathcal D}[f(x)=1]=1/2\pm A^{-1/c_\gamma}$, and,
\item $f$ is resilient: $I_{q,t,\gamma}(f)\le c_\gamma\cdot q/A^{1-\gamma}$.
\end{itemize}
\end{theorem}
